\documentclass[11pt,a4paper]{article}
\usepackage[utf8]{inputenc}
\usepackage[T1]{fontenc}
\usepackage{amsmath, amssymb, amsthm}
\usepackage{xurl} % Better URL breaking
\usepackage{hyperref}
\hypersetup{hidelinks} % Final check to remove all blue boxes
\usepackage{geometry}
\geometry{margin=1in}

\title{\textbf{Spectral Invariants of Star-Cluster Graphs and Topological Regularization of Simplicial Flows}}
\author{Brendan Siche\thanks{The author acknowledges the use of generative AI tools (Claude/Anthropic, Gemini/Google) for numerical verification and typesetting assistance. See Declarations section.}}
\date{March 2026}

\begin{document}

\maketitle

\begin{abstract}
We present current numerical and algebraic investigations into the topological structure of 3D vortex interactions. We identify a scale-invariant constant $R \approx 1.85731$ associated with clustered star graphs and explore its possible role as a dynamic regulator in Navier-Stokes flows. While numerical models in the discrete regime suggest a suppression of enstrophy growth, bridging these findings to the continuous 3D partial differential equations (PDE) remains an open challenge. We document these results as a structured pathway for exploring nonlinear depletion and cascade suppression, rather than as a definitive resolution of the regularity problem.
\end{abstract}

\section{Introduction}
The 3D Navier-Stokes regularity problem remains one of the most profound challenges in mathematical physics. While current analytical bounds allow for potential singularity formation, numerical and experimental evidence often suggest a "depletion" of nonlinearity. We investigate a \textbf{Topological Saturation Hypothesis}, proposing that the geometry of vortex interactions may impose a dynamic limit on stretching efficiency. By mapping these interactions onto spectral graph invariants ($R \approx 1.85731$), we seek to quantify the mechanisms of enstrophy suppression.

\section{Numerical Observations and The Discrete Gap}

\subsection{Discrete Vortex Scaling}
Observations in $N$-point vortex models suggest a logarithmic growth of enstrophy, effectively bounded by the structural constraints of the interaction graph. However, we acknowledge the \textbf{Dimensional Mismatch}: a finite system of ODEs possesses inherent regularity that does not automatically extend to the infinite-dimensional PDE limit. The convergence of consistent discrete solvers ($N^{-0.60}$) provides a robust framework for approximation but does not constitute an analytical proof of global smoothness.

\subsection{Enstrophy Suppression Mechanisms}
We identify three independent mechanisms for the suppression of the forward enstrophy cascade:
\begin{enumerate}
    \item \textbf{Helical De-alignment}: The removal of heterochiral triadic transfers.
    \item \textbf{Coriolis Dispersion}: Wave-mediated suppression in rotating frames.
    \item \textbf{MHD Redistribution}: Magnetic field interactions (Alfven waves) at moderate field strengths.
\end{enumerate}
These mechanisms suggest that the "Monster" of singularity formation is frequently attenuated by physical symmetries and external fields.

\section{Conclusion}
The investigation into the 3D Navier-Stokes equations reveals a complex landscape of dynamic regulators. While the 1.85731 invariant provides a compelling "Geometric Wall" in discrete models, its role in the continuum PDE remains conjectural. We conclude that while the search for regularity remains open, the identification of universal suppression constants offers a promising direction for future research into the stability of incompressible flows.
id.

\section*{Declarations}
\subsection*{Disclosure of Technical Assistance}
In accordance with author guidelines regarding generative tools, the author discloses the use of generative AI tools (Claude/Anthropic and Gemini/Google) for numerical verification, adversarial auditing, and typesetting assistance. All mathematical claims have been independently verified by the author and via reproducible computational scripts. The research design, conjectural framework, and mathematical arguments represent the intellectual work of the human author.

\subsection*{Availability of Code and Data}
The computational results and verification scripts described in this work are available at the following repository:
\begin{center}
    \url{https://github.com/Architect-Resonance/Fluid-Resonance}
\end{center}

\end{document}
